\section{Introduction}\label{sec:intro}
% Numbering map to the project spec (process_geometry_experiment_prompt.md):
%   paper RQ1 = spec RQ1 reframed as a descriptive question ("what is learned"),
%               with no objective-level hypothesis attached;
%   paper RQ2 = spec RQ3, paper RQ3 = spec RQ5;
%   paper H1 = spec H4, H2 = spec H6, H3 = spec H7.
% Spec H1/H2 (objective effects) are reported only as a secondary, Family-B-only
% contrast inside RQ1. Spec RQ2/RQ4 and H3/H5 are deferred to the journal
% extension (see STATUS.md).

Neural sequence models are the default approach to predictive process
monitoring (PPM). They include recurrent~\cite{tax2017lstm,evermann2017deep,camargo2019lstm},
Transformer~\cite{bukhsh2021processtransformer,wuyts2024sutran},
graph~\cite{dissegna2025gnn,wang2025rlhgnn} and fine-tuned language
models~\cite{pasquadibisceglie2024lupin}. They are compared on task metrics
such as accuracy, suffix similarity and remaining-time
error~\cite{rama2022review,teinemaa2019outcome}. These metrics summarise what
a model \emph{outputs}, not what it has \emph{learned}. Every such model maps
each prefix $h_t=(e_1,\dots,e_t)$ to a vector $\z_t\in\R^d$ before its
prediction head, so a case becomes a latent trajectory $\z_1\to\cdots\to\z_T$.
This vector is the model's entire summary of the running case, and its
geometry is never inspected. Outside process mining, such geometry is known to
be informative: perception and language models \emph{straighten} the
trajectories of natural sequences~\cite{henaff2019straightening,hosseini2023straighten}.
In PPM, representations have been treated only as a means to better
features~\cite{dekoninck2018act2vec}.

We study the representations of nine PPM models through three research
questions:
\begin{description}[leftmargin=1.2em,labelsep=0.4em,itemsep=0pt,topsep=2pt,font=\normalfont\bfseries]
  \item[RQ1] What do PPM models learn in their prefix representations? We
        examine how many dimensions they actually use, whether they collapse,
        what shape case trajectories take, and how this varies across
        architectures, objectives and logs.
  \item[RQ2] Does geometry carry information about prediction errors beyond
        confidence and prefix length?
  \item[RQ3] Do models place prefixes with similar \emph{futures} close together
        even when their histories differ?
\end{description}
RQ1 is descriptive. RQ2 matters in practice, because a usable error signal
would tell users when to trust a prediction. RQ3 asks whether a model's state
encodes where a case is going or only where it has been. Answering these
questions requires care: a straight trajectory may simply be a collapsed one,
latent proximity may only reflect shared history, and published models differ
in many respects at once. Our design addresses each of these pitfalls.

\paragraph{Contributions.}
\begin{enumerate*}[label=(\arabic*)]
\item A measurement framework for latent prefix trajectories. It pairs every
      trajectory metric with collapse diagnostics, tests error prediction with
      nested models, and adds two history-controlled tests of \emph{future
      equivalence}. Every metric is validated on synthetic trajectories with
      known values (\cref{sec:method}).
\item A two-family design. Seven published models are reproduced under one
      shared chronological pipeline. A controlled Transformer pair varies only
      the objective, so that any objective effect can be attributed
      (\cref{sec:setup}).
\item An empirical study of RQ1--RQ3 on five public logs, with falsifiable
      hypotheses for RQ2 and RQ3 whose rejections are reported
      (\cref{sec:results}).
\end{enumerate*}
\TODO{One sentence with the headline findings once \cref{sec:results} is filled.}
